\newcommand{\titulus}{\nomenFesti{In Dedicatione Basilicæ S. Mariæ.}
\dies{Die 5. Augusti.}}
\newcommand{\sineabsolutio}{Absolutio de B.M.V.}
\newcommand{\oratio}{\pars{Oratio.}

\noindent Famulórum tuórum, quǽsumus, Dómine, delíctis ignósce, ut, qui tibi placére de áctibus nostris non valémus, Genetrícis Fílii tui intercessióne salvémur.

\pars{Pro pace in universo mundo.} \scriptura{Sir. 50, 25; 2 Esdr. 4, 20; \textbf{H416}}

\vspace{-4mm}

\antiphona{II D}{temporalia/ant-dapacemdomine.gtex}

\vfill

\noindent Deus, a quo sancta desidéria, recta consília et iusta sunt ópera: da servis tuis illam, quam mundus dare non potest, pacem; ut et corda nostra mandátis tuis dédita, et hóstium subláta formídine, témpora sint tua protectióne tranquílla.

\noindent Per Dóminum nostrum Iesum Christum, Fílium tuum, qui tecum vivit et regnat in unitáte Spíritus Sancti, Deus, per ómnia sǽcula sæculórum.

\noindent \Rbardot{} Amen.}
\newcommand{\invitatorium}{\pars{Invitatorium.}

\vspace{-4mm}

\antiphona{V}{temporalia/inv-christummariaefilium.gtex}}
\newcommand{\hymnusmatutinum}{\pars{Hymnus.}

\vspace{-5mm}

\antiphona{II}{temporalia/hym-QuemTerra-alt.gtex}}
\newcommand{\matversus}{\noindent \Vbardot{} María conservábat ómnia verba hæc.

\noindent \Rbardot{} Cónferens in corde suo.}
\newcommand{\lectioi}{\vspace{-2mm}

\sineinitiali{temporalia/oratiodominica-mat.gtex}

\vspace{5mm}

\pars{Absolutio.}

\cuminitiali{}{temporalia/absolutio-precibus.gtex}

\vfill
\pagebreak

\cuminitiali{}{temporalia/benedictio-solemn-noscum.gtex}

\vspace{7mm}

\pars{Lectio I.} \scriptura{2 Cor. 3, 7-18; 4, 1-6}

\noindent De Epístola secúnda beáti Pauli apóstoli ad Corínthios.

\noindent Fratres: Si ministrátio mortis lítteris deformáta in lapídibus fuit in glória, ita ut non possent inténdere fílii Israel in fáciem Móysis propter glóriam vultus eius, quæ evacuátur, quómodo non magis ministrátio Spíritus erit in glória? Nam si ministérium damnatiónis glória est, multo magis abúndat ministérium iustítiæ in glória. Nam nec glorificátum est, quod cláruit in hac parte, propter excelléntem glóriam; si enim, quod evacuátur, per glóriam est, multo magis, quod manet, in glória est.

\noindent Habéntes ígitur talem spem multa fidúcia útimur, et non sicut Móyses: ponébat velámen super fáciem suam, ut non inténderent fílii Israel in finem illíus, quod evacuátur. Sed obtúsi sunt sensus eórum. Usque in hodiérnum enim diem idípsum velámen in lectióne Véteris Testaménti manet non revelátum, quóniam in Christo evacuátur; sed usque in hodiérnum diem, cum légitur Móyses, velámen est pósitum super cor eórum. \emph{Quando autem convérsus fúerit ad Dóminum, aufértur velámen.} Dóminus autem Spíritus est; ubi autem Spíritus Dómini, ibi libértas. Nos vero omnes reveláta fácie glóriam Dómini speculántes, in eándem imáginem transformámur a claritáte in claritátem tamquam a Dómini Spíritu.

\noindent Ideo habéntes hanc ministratiónem, iuxta quod misericórdiam consecúti sumus, non defícimus, sed abdicávimus occúlta dedécoris non ambulántes in astútia neque adulterántes verbum Dei, sed in manifestatióne veritátis commendántes nosmetípsos ad omnem consciéntiam hóminum coram Deo.

\noindent Quod si étiam velátum est evangélium nostrum, in his, qui péreunt, est velátum, in quibus deus huius sǽculi excæcávit mentes infidélium, ut non fúlgeat illuminátio evangélii glóriæ Christi, qui est imágo Dei. Non enim nosmetípsos prædicámus sed Iesum Christum Dóminum; nos autem servos vestros per Iesum. Quóniam Deus qui dixit: «De ténebris lux splendéscat», ipse illúxit in córdibus nostris ad illuminatiónem sciéntiæ claritátis Dei in fácie Iesu Christi.

\noindent \Vbardot{} Tu autem, Dómine, miserére nobis.
\noindent \Rbardot{} Deo grátias.}
\newcommand{\responsoriumi}{\pars{Responsorium 1.} \scriptura{\Rbardot{} Cf. Mt. 17, 2.5.6 \Vbardot{} ibid. 17, 3; \textbf{H160}}

\vspace{-5mm}

\responsorium{VIII}{temporalia/resp-splendidafactaest.gtex}{}}
\newcommand{\lectioii}{\vspace{-7mm}

\cuminitiali{}{temporalia/benedictio-solemn-ipsavirgo.gtex}

\vspace{7mm}

\pars{Lectio II.} \scriptura{Hom. 4: PG 77, 991. 995-996}

\noindent Ex Homília sancti Cyrílli Alexandríni epíscopi in Concílio Ephesíno hábita.

\noindent Sanctórum cœtum, qui a sancta et Deípara sempérque Vírgine María invitáti prompto ánimo huc confluxérunt, lætum erectúmque conspício. Quare licet multa prémerer mæstítia, áttamen hic sanctórum patrum conspéctus lætítiam mihi prǽbuit. Nunc dulce illud hymnógraphi Dávidis verbum apud nos implétum est: \emph{Ecce quid bonum aut quid iucúndum, nisi habitáre fratres in unum?}

\noindent Salve ítaque a nobis, Sancta mýstica Trínitas, quæ nos omnes in hanc Sanctæ Maríæ Deíparæ ecclésiam convocásti.

\noindent Salve a nobis, Deípara María, venerándus totíus orbis thesáurus, lampas inexstinguíbilis, coróna virginitátis, sceptrum rectæ doctrínæ, templum indissolúbile, locus eius qui loco capi non potest, mater et virgo, per quam is benedíctus in sanctis Evangéliis nominátur, \emph{qui venit in nómine Dómini.}

\noindent Salve, quæ imménsum incomprehensúmque in sancto virgíneo útero comprehendísti; per quam Sancta Trínitas glorificátur et adorátur; per quam pretiósa crux celebrátur et in univérso orbe adorátur; per quam cælum exsúltat; per quam ángeli et archángeli lætántur; per quam dǽmones fugántur; per quam tentátor diábolus cælo décidit; per quam prolápsa creatúra in cælum assúmitur; per quam univérsa creatúra, idolórum vesánia deténta, ad veritátis agnitiónem pérvenit; per quam sanctum baptísma obtíngit credéntibus; per quam exsultatiónis óleum, per quam toto terrárum orbe fundátæ sunt Ecclésiæ, per quam gentes adducúntur ad pæniténtiam.

\noindent \Vbardot{} Tu autem, Dómine, miserére nobis.
\noindent \Rbardot{} Deo grátias.}
\newcommand{\responsoriumii}{\pars{Responsorium 2.} \scriptura{\Vbardot{} Lc. 1, 28; \textbf{H47}}

\vspace{-5mm}

\responsorium{II}{temporalia/resp-sanctaetimmaculata-CROCHU.gtex}{}}
\newcommand{\lectioiii}{\vspace{-7mm}

\cuminitiali{}{temporalia/benedictio-solemn-pervirginem.gtex}

\vspace{7mm}

\pars{Lectio III.}

\noindent Et quid plura dicam? per quam unigénitus Dei Fílius iis, \emph{qui in ténebris et in umbra mortis sedébant,} lux resplénduit; per quam prophétæ prænuntiárunt; per quam Apóstoli salútem géntibus prædicárunt; per quam mórtui exsuscitántur; per quam reges regnant, per Sanctam Trinitátem.

\noindent Ecquis hóminum laudabilíssimam illam Maríam pro dignitáte celebráre queat? Ipsa et mater et virgo est; o rem admirándam! Miráculum hoc me in stupórem rapit. Quis umquam audívit ædificatórem prohibéri ne próprium templum, quod ipse constrúxerit, inhabitáret? Quis ob id ignomíniæ sit obnóxius quod própriam fámulam in matrem ascíscat?

\noindent Ecce ígitur ómnia gaudent; contíngat autem nobis ut uniónem revereámur et adorémus, ac indivísam Trinitátem tremámus et colámus, Maríam semper Vírginem, sanctum vidélicet Dei templum, eiusdémque Fílium et sponsum immaculátum láudibus celebrántes: quóniam ipsi glória in sǽcula sæculórum. Amen.

\noindent \Vbardot{} Tu autem, Dómine, miserére nobis.
\noindent \Rbardot{} Deo grátias.}
\newcommand{\responsoriumiii}{\pars{Responsorium 3.} \scriptura{\Vbardot{} Sedulius; \textbf{H48}}

\vspace{-5mm}

\responsorium{VII}{temporalia/resp-congratulaminiquiacum-CROCHU-cumdox.gtex}{}}
\newcommand{\hymnuslaudes}{\pars{Hymnus}

\cuminitiali{I}{temporalia/hym-AveMarisStella.gtex}}
\newcommand{\lectiobrevis}{\pars{Lectio Brevis.} \scriptura{Cf. Is. 61, 10}

\noindent Gaudens gaudébo in Dómino, et exsultábit ánima mea in Deo meo, quia índuit me vestiméntis salútis et induménto iustítiæ circúmdedit me, quasi sponsam ornátam monílibus suis.}
\newcommand{\responsoriumbreve}{\pars{Responsorium breve.} \scriptura{Lc. 1, 28}

\cuminitiali{VI}{temporalia/resp-avemaria-alt.gtex}}
\newcommand{\preces}{\noindent Salvatórem nostrum celebrántes,~\gredagger{} qui ex María Vírgine nasci dignátus est,~\grestar{} exorémus dicéntes:

\Rbardot{} Intercédat pro nobis mater tua, Dómine.

\noindent Salvátor mundi,~\gredagger{} qui redemptiónis tuæ virtúte ab omni peccáti labe matrem tuam præservásti,~\grestar{} serva nos mundos a peccáto.

\Rbardot{} Intercédat pro nobis mater tua, Dómine.

\noindent Redémptor noster,~\gredagger{} qui Vírginem Maríam thálamum puríssimum habitatiónis tuæ et Spíritus Sancti fecísti sacrárium,~\grestar{} nos templum tui Spíritus fac perénne.

\Rbardot{} Intercédat pro nobis mater tua, Dómine.

\noindent Verbum ætérnum, quod matrem tuam docuísti óptimam sibi partem elígere,~\grestar{} tríbue nobis eam imitári, cibum quæréntes, qui permáneat in vitam ætérnam.

\Rbardot{} Intercédat pro nobis mater tua, Dómine.

\noindent Rex regum, qui matrem tuam córpore et ánima tecum voluísti in cælum assúmptam,~\grestar{} fac ut quæ sursum sunt semper cogitémus.

\Rbardot{} Intercédat pro nobis mater tua, Dómine.

\noindent Dómine cæli et terræ, qui Maríam regínam a dextris tuis astáre fecísti,~\grestar{} tríbue nos eiúsdem glóriæ meréri consórtium.

\Rbardot{} Intercédat pro nobis mater tua, Dómine.}
\newcommand{\benedictus}{\pars{Canticum Zachariæ.}

\vspace{-4mm}

\antiphona{VI C}{temporalia/ant-mentesnostrascastifica.gtex}

\vspace{-2mm}

\scriptura{Lc. 1, 68-79}

\vspace{-2mm}

\cantusSineNeumas
\initiumpsalmi{temporalia/benedictus-initium-vi-C-auto.gtex}

%\vspace{-1.5mm}

\input{temporalia/benedictus-vi-C.tex} \Abardot{}}
\newcommand{\benedicamuslaudes}{\cuminitiali{VIII}{temporalia/benedicamus-officium-bmv.gtex}}
\include{hebdomadaxviii}
\include{feriaiv}
